\section{Related literature}
\label{sec:literature}

This reassessment sits at the intersection of switching-cost competition and behavior- or history-based pricing. Switching costs make inherited market shares strategically valuable because firms trade off harvesting attached consumers against competing for rivals' consumers; \citet{Klemperer1995} surveys this broad mechanism. In dynamic models, customer acquisition affects future states and prices. Examples include \citet{Chen1997}, \citet{VillasBoas1999}, and, in a continuous-time setting with asymmetric market shares, \citet{FabraGarcia2015}. Our result is not a dynamic switching-cost theorem: inherited shares are fixed primitives of the one-period pricing game considered by \citet{GehrigShyStenbacka2012}.

A second strand studies customer recognition and targeted poaching. \citet{FudenbergTirole2000} analyze customer poaching and brand switching under alternative contract structures, while \citet{ShafferZhang2000} study preference-based discrimination with switching costs. \citet{AcquistiVarian2005} examine conditioning prices on purchase histories, including competitive environments. These papers establish that customer information can intensify competition and alter switching, but they do not supply the global best-response condition in the particular inherited-history Hotelling benchmark studied here.

A particularly close recent neighborhood combines asymmetric firms, switching costs, and behavior-based pricing. \citet{Umezawa2022} studies a two-period horizontally and vertically differentiated duopoly. An author-issued corrigendum, \citet{UmezawaCorrigendum2023}, is methodologically relevant here: it corrects the uniform-pricing benchmark because a price ordering had been used to infer a switching/allocation ordering that need not follow. The corrected analysis classifies regimes by the actual allocation cutoffs instead. This is a direct precedent for treating regime consistency as a mathematical obligation rather than as a consequence of a candidate price ordering.

The Umezawa model is nevertheless not identified here as a parent theorem for the present correction. Its timing is two-period, horizontal and vertical differentiation jointly determine demand, and customer histories are generated dynamically. The present paper instead takes the GSS inherited split as a primitive and solves a one-period uniform-price game. Likewise, \citet{Shrivastav2023} changes customer recognition through imperfect information and misinformation; \citet{ColomboGrazianoPignataro2024} studies imperfect history-based recognition with asymmetric market shares; and \citet{UmezawaYamakawa2025} introduces multiple consumer types and a different switching-cost incidence. We did not establish a formal nesting or non-nesting theorem for these models. The narrower claim is that our search located no prior derivation of the GSS-specific threshold \(x_H(s)\) or its complete pure-price correspondence.

The closest policy-choice comparison is \citet{JeongMaruyama2008}, where firms choose between uniform and discriminatory pricing under exogenous switching costs. Their equilibrium object includes the pricing-policy choice itself, whereas our object is the complete pure-strategy correspondence of the fixed uniform-price subgame in the GSS benchmark.

Accordingly, the contribution is narrower than a new general theory of behavior-based discrimination. It is a mathematical reassessment of a specific benchmark that is frequently interpretable within this literature: the accepted uniform-price solution is a local switching-branch candidate, and its global pure-equilibrium validity requires comparing it with a no-poaching kink. The resulting threshold changes which downstream welfare comparisons are equilibrium comparisons at all.
